\subsection{Choosing a calibration scale}

To resolve reference disagreement at a scale of one in a thousand,
the independent-slot model in \cref{sec:6-2} suggests
\[
 \frac{4.74}{Sk}\simeq10^{-3},\quad k=10
 \quad\Longrightarrow\quad S\simeq474\approx500.
\]
This motivates testing the whole 500-query calibration pool alongside
smaller samples. The model concerns reference agreement; the measured
held-out quality and sampled risks are reported in
\cref{tab:7-budget-full,fig:7-budget-risk}. The search-order conditions
for full-pool bisection remain those in \cref{sec:6-3}.

\subsection{Risk and selected-budget performance}

Let $m$ random calibration samples select budgets $\widehat\alpha_r$.
Using the best tested held-out recall as the reference, define the
loss for sample $r$ and the fraction whose loss exceeds the target:
\begin{equation}
 L_r=\widehat r^*-\widehat r(\widehat\alpha_r),\qquad
 \widehat p_\delta=\frac{1}{m}\sum_{r=1}^{m}\mathbf1\{L_r>\delta\}.
 \label{eq:6-risk}
\end{equation}
Mean loss describes average quality, the frequency of exceeding the target describes how
often the target is missed, and tail loss describes how large the misses
are. \Cref{tab:7-budget-quality,fig:7-budget-risk} reports these summaries;
the largest per-sample losses are in \cref{sec:7-budget-trace}.

The reported rate
$\overline Q=m^{-1}\sum_r\widehat Q(\widehat\alpha_r)$ averages the
previously measured throughput of the budget each sample chose.
It is an estimate from fixed-budget runs, distinct from the rate at the
median budget and from directly timing a calibrated query stream.
\Cref{tab:7-budget-cost} reports it alongside quality.

\subsection{Cost of reuse}

Calibration pays for itself only if later searches save enough time.
Let $T_{\rm cal}$ be its one-time cost, $T_{\rm chosen}$ the time of a
batch at the chosen budget and $T_f$ that batch's time at a fixed budget.
For $B_{\rm reuse}$ equal-sized subsequent batches,
\begin{equation}
 T_{\rm policy}(B_{\rm reuse})=T_{\rm cal}+B_{\rm reuse}T_{\rm chosen},\qquad
 T_{\rm fixed}(B_{\rm reuse})=B_{\rm reuse}T_f.
 \label{eq:6-total-cost}
\end{equation}
When $T_{\rm chosen}<T_f$, dividing the setup cost by the saving per
batch gives the break-even count:
\begin{equation}
 B_{\rm reuse}^*=\frac{T_{\rm cal}}{T_f-T_{\rm chosen}}.
 \label{eq:6-payback}
\end{equation}
The strict time-saving condition follows directly:
\[
 \begin{aligned}
 T_{\rm policy}<T_{\rm fixed}
 &\iff T_{\rm cal}<B_{\rm reuse}(T_f-T_{\rm chosen})\\
 &\iff B_{\rm reuse}>B_{\rm reuse}^*
 \qquad(T_{\rm chosen}<T_f).
 \end{aligned}
\]
If $T_{\rm chosen}\geq T_f$, no number of batches repays calibration
through throughput. For $n$ queries per batch,
$T=1000n/Q$ converts QPS to milliseconds. When chosen budgets vary,
$1/\mathbb E[Q]\leq\mathbb E[1/Q]$: inverting mean QPS understates
mean batch time. The timing boundaries and direct stream checks are in
\cref{sec:a-budget-contract,sec:7-stream-cost}; model estimates are in
\cref{tab:7-budget-cost}.

Recalibration every $H$ batches adds $T_{\rm cal}/H$ per batch.
Automatic detection of workload changes and renewal of calibration
remain future work.
